\vspace{-0.3in}
\section{Related Work}
\vspace{-0.14in}
Manufacturer-specified operational parameters for general deployment of robots are typically far more conservative than hardware-level physical limits and often lack transparency in their derivation.
This has led to significant under-utilization of robotic capabilities, as evidenced by geometric motion planning and time parameterization approaches dominating most industrial use \cite{berscheid2021jerk,coleman2014reducing,orthey2023sampling} and recent advances in foundation models for manipulation in unstructured environments being restricted to exclusively geometric reasoning tasks \cite{o2024open,team2024octo,zitkovich2023rt}.
A nuanced, configuration-dependent understanding of the robot’s specifications that considers both robot and environment dynamics could enable a broader spectrum of manipulation capabilities within hardware limits.

Higher-order constraints such as joint torques and end-effector accelerations can be addressed through several approaches that range from ``plan-and-filter'' methods which post-process geometric paths \cite{abderezaei2024clutter,kuntz2019fast} to kinodynamic planners that directly incorporate dynamics during planning \cite{li2015sparse,nayak2022bidirectional,pasricha2024virtues}.
While these approaches offer completeness guarantees, they face practical limitations in high-dimensional manipulation spaces and often exhibit high variance in planning times especially when faced with state or model uncertainty.
These challenges are particularly acute when planning near operational limits, where the feasible solution space becomes increasingly constrained \cite{kingston2018sampling}.
Optimization-based methods naturally accommodate dynamics constraints but rely heavily on good initialization, leading to hybrid approaches that bootstrap the optimization process with learning or sampling-based solutions \cite{ichnowski2020deep,kuntz2019fast,yan2024impact}.
Optimal control formulations for dynamics-aware manipulation also exist \cite{arrizabalaga2024slosh,kim2014catching}. Despite their theoretical rigor, they face practical limitations in handling complex constraints and achieving reliable convergence for real-world manipulation problems.

Importantly, there exist a large body of works that demonstrate how a robot's behavioral repertoire can be significantly expanded when incorporating dynamic constraints---e.g. nonprehensile manipulation \cite{dogar2011framework,pasricha2022pokerrt,ruggiero2018nonprehensile,yu2016more,zeng2020tossingbot} or fast inertial transport \cite{arrizabalaga2024slosh,ichnowski2022gomp,muchacho2022solution,avigal2022gomp}.
These capabilities and considerations are particularly relevant for the specific problem of payload manipulation, where the interaction between payload and induced joint torques significantly influence task success.
Interestingly, payload manipulation presents unique challenges that intersect robot design and control.
Traditionally, prior work has focused on mechanical modeling for specific low-DoF systems \cite{wang2001payload}, with emphasis on payload capacity optimization \cite{korayem2008maximum} and joint torque considerations \cite{nho2003intelligent}. Other work in this domain tackles super-nominal payload handling through the use of multi-arm systems \cite{kim2021enhancing,yan2016coordinated}.
However, the field lacks generalizable approaches for high-DoF systems that can reason about dynamic payload-robot interactions across diverse tasks.

Learning-based approaches offer promising directions for both generalizability and efficient trajectory generation in high-dimensional constrained spaces while respecting system dynamics. Autoregressive models implicitly encode dynamics through dataset design \cite{fishman2023motion,qureshi2020motion}, while recent diffusion-based approaches provide non-autoregressive alternatives with improved sampling diversity \cite{carvalho2023motion,saha2024edmp}. Task-level diffusion policies, conditioned on language and visual embeddings \cite{ajay2023is,chi2023diffusionpolicy}, demonstrate potential for generating diverse behaviors from high-level task specifications.
Our method conditions a denoising diffusion model on the target payload to generate dynamically feasible trajectories.
